% AY2022 材料力学 〔II〕（転記）
% 出典: 04_材料力学/original/04_材料力学/2022/2022za.pdf
% 要約: 長さLの丸棒(直径d)は点Aで固定, 他端Bに長さL/3の剛体棒が接合。丸棒先端Bに下向き
%       集中荷重P, 同時に剛体棒先端Cにも下向き集中荷重P。縦弾性係数E, 横弾性係数G。
%       (1) 固定端A支持力Ra・支持モーメントMa・支持トルクTa, (2) SFD/BMD,
%       (3) AからL/4位置断面の点uのσx,σy,τxy, (4) B,C点のz方向変位δB,δC（断面二次I, 極Ip）。
% rem: 原本は3次元俯瞰図。丸棒(x方向)+剛体棒(y方向)のL字構成を簡略立体図で再現する。

\section*{〔II〕}

図2のように, 長さ $L$ の丸棒（直径 $d$）は点 A で固定され, 他端 B に長さ $L/3$ の剛体棒が
接合されている. 丸棒の先端 B に下向きに集中荷重 $P$ が作用し, 同時に剛体棒の先端 C にも
下向きに集中荷重 $P$ が作用する場合について, 問い (1)\,$\sim$\,(4) に答えよ.
丸棒の縦弾性係数を $E$, 横弾性係数を $G$ とする.

\begin{center}
\begin{tikzpicture}[>=Latex]
  % 固定壁（A, z-up 斜め平面）
  \fill[pattern=north east lines] (0.2,0.6) -- (0.2,3.0) -- (1.2,3.4) -- (1.2,1.0) -- cycle;
  \draw (0.2,0.6) -- (0.2,3.0) -- (1.2,3.4) -- (1.2,1.0) -- cycle;
  % 丸棒 A → B（x方向, 下右へ）
  \draw[line width=12pt,gray!25,line cap=round] (0.9,2.4) -- (6.0,1.0);
  \draw[fill=gray!25] (6.0,1.0) ellipse (0.32 and 0.12);
  % 剛体棒 B → C（y方向, 右へ）
  \draw[line width=2pt] (6.0,0.95) -- (8.3,1.5);
  % 軸（A点）
  \draw[->] (0.9,2.4) -- (0.9,3.7) node[left]{$z$};
  \draw[->] (0.9,2.4) -- (2.0,2.95) node[above]{$y$};
  \draw[->] (0.9,2.4) -- (2.0,1.95) node[below]{$x$};
  % 節点
  \node[left] at (0.8,2.35) {A}; \node[below] at (6.0,0.75) {B}; \node[right] at (8.35,1.5) {C};
  % 点 u（L/4, 上面）
  \fill (2.3,2.25) circle (1.8pt); \node[above] at (2.3,2.3) {u};
  % 荷重 P（B, C 下向き）
  \draw[->,very thick] (6.0,2.1) -- (6.0,1.15); \node[right] at (6.05,1.95) {$P$};
  \draw[->,very thick] (8.3,2.55) -- (8.3,1.6); \node[right] at (8.35,2.4) {$P$};
  % 寸法 L（A-B）, L/4（A-u）, L/3（B-C）
  \draw[<->] (1.1,3.05) -- (6.0,1.65) node[midway,above,font=\scriptsize]{$L$};
  \draw[<->] (1.0,1.95) -- (2.3,1.6) node[midway,below,font=\scriptsize]{$L/4$};
  \draw[<->] (6.1,0.65) -- (8.3,1.2) node[midway,below,font=\scriptsize]{$L/3$};
  % u断面インセット（左下）
  \begin{scope}[shift={(0.9,-1.5)}]
    \draw (0,0) circle (0.55);
    \draw[->] (0,-0.85) -- (0,0.95) node[right]{$z$};
    \draw[->] (-0.85,0) -- (0.95,0) node[above]{$y$};
    \fill (0,0.55) circle (1.8pt); \node[above left] at (0,0.55) {u};
    \node at (0,-1.25) {\underline{断面}};
  \end{scope}
\end{tikzpicture}

\vspace{1mm}
図2
\end{center}

\begin{enumerate}[label=(\arabic*)]
  \item 固定端 A での支持力 $R_a$, 支持モーメント $M_a$, 支持トルク $T_a$ を図示し, 数式で示せ.

  \item 丸棒のせん断力図（SFD）と曲げモーメント図（BMD）を求めて, 図示せよ.

  \item 固定端 A から $L/4$ の位置にある断面の点 u の $x$ 方向の応力 $\sigma_x$,
        $y$ 方向の応力 $\sigma_y$, せん断応力 $\tau_{xy}$ を求めよ.

  \item 点 B および C の $z$ 方向の変位 $\delta_B$, $\delta_C$ を求めよ. ただし, 丸棒の
        断面二次モーメントは $I$, 断面二次極モーメントは $I_p$ とする.
\end{enumerate}
